\chapter{Cinemática y control de la base móvil}
\label{cap:movilidad}

\section{Referencias, unidades y geometría}
Se define un marco de cuerpo $B$ con $x$ hacia adelante, $y$ a la izquierda y $z$ hacia arriba. El giro positivo $\psi$ es antihorario visto desde arriba. Las velocidades de cuerpo son
\begin{equation}
\vct{V}_B=\begin{bmatrix}v_x&v_y&\omega\end{bmatrix}^{\trans},
\qquad v_x,v_y\ [\mathrm{m/s}],\quad\omega\ [\mathrm{rad/s}].
\end{equation}
El orden lógico de ruedas es delantera izquierda (FL), delantera derecha (FR), trasera izquierda (RL) y trasera derecha (RR). Se usa radio nominal $r=0{,}040\unit m$ y suma de semidistancias $L=l_x+l_y=0{,}170\unit m$. El contacto efectivo de rodillos puede diferir de esa geometría nominal.

Para una rueda situada en $(x_i,y_i)$, la velocidad del punto de montaje inducida por un movimiento rígido es
\begin{equation}
\vct{v}_i=\begin{bmatrix}v_x-\omega y_i\\v_y+\omega x_i\end{bmatrix}.
\end{equation}
La proyección según los rodillos a 45 grados produce las combinaciones de avance, lateral y giro. La disposición física de rodillos debe corresponder a la matriz utilizada. El modelo sigue el formalismo de robots móviles omnidireccionales~\cite{omni}.

\section{Cinemática inversa Mecanum}
Sea $\vct{u}=[u_{\mathrm{FL}},u_{\mathrm{FR}},u_{\mathrm{RL}},u_{\mathrm{RR}}]^{\trans}$ el vector de velocidades periféricas de rueda. El mezclador geométrico es
\begin{equation}
\vct{u}=\mat{A}\mat{S}\vct{V}_B,
\qquad
\mat{A}=\begin{bmatrix}
1&-1&-L\\1&1&L\\1&1&-L\\1&-1&L
\end{bmatrix}.
\label{eq:mecanum}
\end{equation}
La matriz $\mat S$ adapta los sentidos del contrato lógico a los de movimiento. Para la configuración descrita:
\begin{equation}
\mat{S}=\diag(-1,-1,-1)=-\mat{I}_3.
\end{equation}
Estos signos no sustituyen las calibraciones eléctricas individuales del motor y del encoder. Primero se normalizan las cuentas y las salidas por rueda; luego se usan las matrices de cuerpo.

La velocidad angular accionada se obtiene como
\begin{equation}
\vct{\omega}_w=\frac{1}{r}\mat A\mat S\vct V_B.
\end{equation}
La distinción es esencial: $u_i$ es velocidad periférica, $\omega_{w,i}$ es velocidad angular y PWM es una acción eléctrica. No son magnitudes intercambiables.

\subsection{Ejemplo de avance}
Para $\vct V_B=[0{,}100\;0\;0]^{\trans}$, en unidades SI, se obtiene
\begin{equation}
\vct u=\begin{bmatrix}-0{,}100\\-0{,}100\\-0{,}100\\-0{,}100\end{bmatrix}\unit{m/s},
\qquad
\vct\omega_w=\begin{bmatrix}-2{,}5\\-2{,}5\\-2{,}5\\-2{,}5\end{bmatrix}\unit{rad/s}.
\end{equation}
La magnitud equivale a aproximadamente $23{,}87\unit{rpm}$. El signo negativo resulta de la base lógica elegida; la polaridad efectiva se aplica en el controlador de cada rueda.

\subsection{Límites y saturación}
Para un límite periférico común $u_{\max}$, un escalado uniforme conserva la dirección solicitada:
\begin{equation}
\gamma=\max\left(1,\max_i\frac{|u_i|}{u_{\max}}\right),
\qquad \vct u_{\mathrm{apl}}=\frac{\vct u}{\gamma}.
\end{equation}
Recortar las ruedas independientemente puede cambiar las proporciones de avance, lateral y giro. El parámetro nominal de 150 mm/s del software es una consigna límite de configuración, sin equivaler a una velocidad máxima física caracterizada.

\section{Cinemática directa y redundancia}
La reconstrucción ideal se obtiene por pseudoinversa. Como $L\ne0$:
\begin{equation}
\mat A^{\trans}\mat A=\diag(4,4,4L^2),
\qquad \mat A^+=(\mat A^{\trans}\mat A)^{-1}\mat A^{\trans}.
\end{equation}
Por tanto,
\begin{equation}
\mat A^+=\frac14\begin{bmatrix}
1&1&1&1\\-1&1&1&-1\\-1/L&1/L&-1/L&1/L
\end{bmatrix},
\qquad \vct V_B=\mat S^{-1}\mat A^+\vct u.
\label{eq:directam}
\end{equation}
En esta configuración $\mat S^{-1}=\mat S$. Cuatro observaciones de rueda determinan tres velocidades de cuerpo bajo el modelo ideal. El residuo
\begin{equation}
\vct\varepsilon=\vct u-\mat A\mat S\widehat{\vct V}_B
\end{equation}
permite observar incompatibilidades entre medidas y modelo. Un deslizamiento coherente de varias ruedas puede no producir un residuo grande, de modo que no reemplaza una referencia externa.

\section{Encoders y conversión de cuentas}
La configuración nominal utiliza $N_i=370$ cuentas por vuelta de rueda y captura x1: flanco ascendente del canal A y lectura de B para el sentido. Ese valor no se multiplica nuevamente por cuatro.

Para un incremento normalizado $\Delta n_i$ y período $\Delta t$:
\begin{align}
\Delta\phi_i&=\frac{2\pi\Delta n_i}{N_i},\\
\Delta s_i&=r_i\Delta\phi_i=\frac{\pi D_i\Delta n_i}{N_i},\\
\widehat u_i&=\frac{\Delta s_i}{\Delta t}.
\end{align}
Con $D=80\unit{mm}$, la cuantización lineal nominal es
\begin{equation}
q_s=\frac{\pi D}{N}=\frac{80\pi}{370}\simeq0{,}679\unit{mm/cuenta}.
\end{equation}
La cuantización no representa la exactitud del desplazamiento: diámetro efectivo, deformación, carga, piso y patinamiento pueden dominar el error. Una ausencia de cuentas puede indicar reposo, falta de conexión o fallo de adquisición; se necesita estado de validez para interpretarla.

\section{Odometría y escalas direccionales}
La inversa Mecanum también se aplica a incrementos de recorrido:
\begin{equation}
\Delta\vct\xi_B=\begin{bmatrix}\Delta x_B\\\Delta y_B\\\Delta\psi\end{bmatrix}
=\mat S^{-1}\mat A^+\begin{bmatrix}\Delta s_{\mathrm{FL}}\\\Delta s_{\mathrm{FR}}\\\Delta s_{\mathrm{RL}}\\\Delta s_{\mathrm{RR}}\end{bmatrix}.
\end{equation}
Se aplican factores empíricos por componente y por sentido antes de transformar al mundo:
\begin{equation}
\Delta\vct\xi_{B,c}=
\underbrace{\begin{bmatrix}k_x&0&0\\0&k_y&0\\0&0&k_\psi\end{bmatrix}}_{\mat K_d}
\Delta\vct\xi_B.
\end{equation}
El signo de cada incremento elige avance o retroceso, izquierda o derecha, y CW o CCW. Los valores de fábrica considerados son 1,000 para avance y retroceso, 0,775 para laterales y 1,028 para giros. Una configuración guardada puede contener otros valores. Estas escalas corrigen odometría; no multiplican PWM ni la meta de misión.

\subsection{Interpretación de la calibración}
Si una misión se detiene al alcanzar una meta odométrica $d_{\mathrm{obj}}$ y se recorre una distancia real $d_{\mathrm{real}}$, el ajuste es
\begin{equation}
k_{\mathrm{nuevo}}=k_{\mathrm{actual}}
\frac{d_{\mathrm{real}}}{d_{\mathrm{obj}}}.
\label{eq:ajustek}
\end{equation}
Si al contar 1000 mm laterales se recorren 775 mm reales, el factor pasa de 1 a 0,775. Para el mismo incremento de cuentas se publica menos desplazamiento y la misión requiere más cuentas para completar la meta. Si la parada la define una distancia física externa, se compara la distancia real con la odometría obtenida sobre ese mismo tramo; no se intercambian ambos protocolos.

Los factores tienen resolución de una milésima, rango 0 a 2, y cero bloquea la componente correspondiente en control cerrado. Modificar la escala exige reiniciar el marco de pose para no mezclar una trayectoria con dos calibraciones.

\subsection{Integración en el mundo}
La relación diferencial es
\begin{equation}
\begin{bmatrix}\dot x\\\dot y\\\dot\psi\end{bmatrix}=
\begin{bmatrix}\cos\psi&-\sin\psi&0\\\sin\psi&\cos\psi&0\\0&0&1\end{bmatrix}
\begin{bmatrix}v_x\\v_y\\\omega\end{bmatrix}.
\end{equation}
La integración de punto medio usa los incrementos corregidos y el rumbo medio
\begin{align}
\psi_m&=\psi_k+\frac{\Delta\psi_c}{2},\\
x_{k+1}&=x_k+\Delta x_{B,c}\cos\psi_m-\Delta y_{B,c}\sin\psi_m,\\
y_{k+1}&=y_k+\Delta x_{B,c}\sin\psi_m+\Delta y_{B,c}\cos\psi_m,\\
\psi_{k+1}&=\wrap(\psi_k+\Delta\psi_c).
\end{align}
La pose es relativa al origen de inicialización. Sin referencia externa acumula deriva y no constituye localización absoluta.

\section{Regulación de velocidad por rueda}
El control compara consigna y velocidad Hall con acumuladores independientes. La implementación usa ganancias comunes en formato Q8 y recortes numéricos:
\begin{align}
e_i[k]&=\sat_{2000}\!\left(u_{\mathrm{ref},i}[k]-\widehat u_i[k]\right),\\
E_i[k]&=\sat_{32767}\!\left(E_i[k-1]+e_i[k]\right),\\
p_i^*[k]&=\sat_{p_{\max}}\!\left(
\frac{K_pe_i[k]+K_iE_i[k]+K_d\bigl(e_i[k]-e_i[k-1]\bigr)}{256}\right).
\end{align}
Las velocidades de esta implementación se expresan en mm/s. Los valores iniciales son $K_p=128$, $K_i=8$ y $K_d=0$. El acumulador suma por ciclo sin multiplicación explícita por tiempo; sus ganancias dependen del período de servicio. Antes de escribir al puente H se aplica una rampa a la salida.

El recorte del acumulador evita desbordamiento, pero no equivale a un antiwindup completo ligado a saturación. La compensación de zona muerta y el ajuste de ganancias requieren medidas. PWM manual y prueba de rueda evitan este lazo, por lo que no tienen las mismas condiciones que una misión cerrada.

\section{Navegación a una pose planar}
Para una meta $(x_d,y_d,\psi_d)$ se calculan errores de mundo y se rotan al cuerpo:
\begin{equation}
\begin{bmatrix}e_{x,B}\\e_{y,B}\end{bmatrix}
=\begin{bmatrix}\cos\psi&\sin\psi\\-\sin\psi&\cos\psi\end{bmatrix}
\begin{bmatrix}x_d-x\\y_d-y\end{bmatrix},
\qquad e_\psi=\wrap(\psi_d-\psi).
\end{equation}
La ley nominal previa a límites es
\begin{equation}
v_x=K_{xy}e_{x,B},\qquad v_y=K_{xy}e_{y,B},\qquad
\omega=K_\psi e_\psi.
\end{equation}
Se utiliza $K_{xy}\simeq1\unit{s^{-1}}$ y ganancia de rumbo de $3\unit{s^{-1}}$ durante traslación, frente a $1\unit{s^{-1}}$ en giro puro. La implementación agrega límites, mínimos para superar zona muerta y normalización de mezcla. El giro puro anula las consignas de traslación.

La configuración de llegada de la fuente 1.28.1 usa
\begin{equation}
|e_x|\le4\unit{mm},\qquad |e_y|\le4\unit{mm},\qquad
|e_\psi|\le10\unit{mrad}\simeq0{,}57^{\circ}.
\end{equation}
Estas bandas son decisiones de software, no una precisión física demostrada. Los resultados experimentales del capítulo~\ref{cap:resultados} corresponden a 1.28.0, con bandas de 10 mm y 25 mrad.
\FloatBarrier
